\scanpage{25}
\margin{$\unsure{0)}$}
\margin{$\big|$}
\uline{Proposition} Soit $x\in\widetilde A$ tel que \inserted{\deleted{\unsure{avec}} $\varepsilon(x)=n$ et} $\lambda^i(x)=0$ pour $i>n$, i.e. tel que $\lambda_t(x)$ soit un polynôme de degré $\unsure{\leq}n$. Alors \deleted{$\unsure{x^i=0}$} $(i-1)!x^i=0$ pour $i>n$.

En effet, $\lambda^n(x)$ est de la forme $[1,1+x^1\mathbin{\text{---}}]$ donc est inversible, par suite $x$ est un élément \uline{géométrique} \inserted{de rg $n$} de $\widetilde A$, donc
\[
\gamma^m(x-\varepsilon(x))=0\quad\text{pour }m>n.
\]
D'où la conclusion.

\uline{Remarque} Pour que $x=[n,1+\text{---}+x^n+x^{n+1}]$ soit géométrique, quand $A^{\unsure{n+2}}=0$, il faut et il \uline{suffit} que: $n!x^{n+1}=0$.
/ Il faut et suffit, posant $X=[n,1+x^1\mathbin{\text{---}}+x^n]$, $Y=[n+1,1+x^{n+1}]$, d'où $x=X+Y-(n+1)$, que $\lambda_t(X)\lambda_t(Y)$ soit divisible par $(1+t)^{n+1}$,
i.e.
\[
D_t^i(\lambda_t(X)\lambda_t(Y))\big|_{t=-1}=0\quad\text{pour }0\leq i\leq n.
\]

\margin{$\Big|$}
\uline{Corollaire} Soit $x$ un élément \deleted{géométrique} d'un $\lambda$-anneau spécial augmenté, \deleted{\unsure{supposons} $\unsure{\varepsilon(x)}$ $\lambda^i(x)$}\inserted{Soit $n=\varepsilon(x)$, et $m>n$}. Alors $\gamma^m(x\deleted{\unsure{-\varepsilon(x)}})$ appartient à \unsure{l'-$\lambda$-idéal} engendré par les $\lambda^i(x)$ ($i>n$) \deleted{et \unsure{le} \deleted{d'un}\inserted{des}}
\deleted{$P((\gamma^i(x-\unsure{n}))\unsure{\gamma^j(\ )})$},
\inserted{$P$ \unsure{isobare} de poids \uuline{$m+1$} en les $\gamma^i(x-n)$}

\deleted{
On peut supposer
\[
A=\mathbb Z[(\lambda^i(X))_{\unsure{i\geq m|i\geq1}}],\qquad
\deleted{\varepsilon(X)=n}\inserted{\text{\unsure{pour} \deleted{\unsure{x}} $X+n$ ($\varepsilon(X)=0$)}}.
\]
\deleted{L'assertion} Soit $I$ l'idéal engendré par les $\lambda^i(X)$ ($i>n$), l'assertion signifie que $\gamma^m(X-n)\deleted{\unsure{\bmod}}I$ \unsure{est dans} l'idéal engendré par les $P(\gamma^i(X-n))$, $P$ isobare de poids $m+1$. Cela signifie la $=$ \unsure{relation} dans
\[
B=\mathbb Z[(\lambda^i(X))_{1\leq i\leq n}]:
\]
Dans cet anneau, $\gamma^m(X-n)$ est de \deleted{``filtration'' $\geq m+1$ (pour $m>n$)}
}
\margin{\deleted{$\unsure{\lambda^i(x-y)}$; $\unsure{\gamma^m(x)}$; $\unsure{\gamma^i(x)\gamma^j(y)}$}}
\margin{24}
